\chapter{\texorpdfstring{$R^2$}{R-squared}, and the trap inside it}
\label{ch:metrics}

Every headline number in the paper is an \Rsq{} value. Almost every subtlety in
this project comes from what those values are computed \emph{over}. This
chapter is therefore not optional background --- it is the key to the whole
thing.

\section{What \texorpdfstring{$R^2$}{R-squared} means}

\Rsq{} answers one question: \emph{how much better are these predictions than
just guessing the average?}

Given true values $y_1,\dots,y_n$ and predictions $\hat{y}_1,\dots,\hat{y}_n$,
let $\bar{y}$ be the mean of the true values. Define

\[
  \mathrm{SS}_{\text{res}} = \sum_{i=1}^{n} (y_i - \hat{y}_i)^2,
  \qquad
  \mathrm{SS}_{\text{tot}} = \sum_{i=1}^{n} (y_i - \bar{y})^2,
\]

\[
  \boxed{\;R^2 = 1 - \frac{\mathrm{SS}_{\text{res}}}{\mathrm{SS}_{\text{tot}}}\;}
\]

$\mathrm{SS}_{\text{res}}$ is how wrong the predictions are.
$\mathrm{SS}_{\text{tot}}$ is how wrong you would be if you always guessed the
mean. The ratio compares the two.

\begin{keybox}{Reading the scale}
\begin{description}[leftmargin=1.6cm,style=nextline,itemsep=1pt]
\item[$R^2 = 1$] Perfect. Every prediction exactly right.
\item[$R^2 = 0$] Exactly as good as always guessing the mean.
\item[$R^2 < 0$] \emph{Worse} than guessing the mean. Entirely possible, and it
  happens frequently in this project.
\end{description}
The negative range surprises people who have only seen \Rsq{} as ``the fraction
of variance explained''. That interpretation holds only for a fitted linear
regression on its own training data. For predictions from anywhere else --- a
neural network, another dataset, a wish list --- there is no lower bound.
\end{keybox}

\section{The trap: \texorpdfstring{$R^2$}{R-squared} over \emph{what}?}
\label{sec:r2trap}

The formula needs a set of pairs. Nothing says what the set is. The paper and
a casual reader can easily mean different things, and here they do.

\begin{figure}[htbp]
\centering
\begin{tikzpicture}[font=\small,
  hdr/.style={font=\small\bfseries},
  cell/.style={draw=greyc!50,minimum width=1.05cm,minimum height=0.52cm,
               font=\scriptsize,inner sep=1pt},
]
% ---------- LEFT: across-dataset ----------
\node[hdr,text=greyc] at (0,3.3) {What most readers assume};
\foreach \i/\lab in {0/seq 1,1/seq 2,2/seq 3,3/seq 4,4/$\vdots$} {
  \node[cell,fill=boxbg] at (-1.6, 2.6-0.55*\i) {\lab};
  \node[cell] at (-0.5, 2.6-0.55*\i) {true};
  \node[cell] at (0.6, 2.6-0.55*\i) {pred};
}
\draw[ourscol,line width=1.2pt,rounded corners=2pt]
  (-1.05,2.88) rectangle (1.15,0.32);
\node[font=\scriptsize,text=ourscol,align=center,text width=3.2cm]
  at (0.05,-0.15)
  {$R^2$ \emph{down} the column:\\ many sequences, one property};

% ---------- RIGHT: within-vector ----------
\node[hdr,text=greyc] at (6.4,3.3) {What the paper computes};
\foreach \i/\lab in {0/tough,1/tough SD,2/$E$,3/$E$ SD,4/$\vdots$} {
  \node[cell,fill=boxbg] at (4.6, 2.6-0.55*\i) {\lab};
  \node[cell] at (5.9, 2.6-0.55*\i) {target};
  \node[cell] at (7.0, 2.6-0.55*\i) {pred};
}
\draw[papercol,line width=1.2pt,rounded corners=2pt]
  (5.35,2.88) rectangle (7.55,0.32);
\node[font=\scriptsize,text=papercol,align=center,text width=3.6cm]
  at (6.45,-0.15)
  {$R^2$ \emph{within one row}:\\ one sequence, eight properties};
\end{tikzpicture}
\caption[Two meanings of R-squared]{The same formula, two completely different
quantities. Left: the standard regression sense --- how well does the model
predict toughness across many sequences? Right: what the paper actually
computes --- for one designed sequence, do its eight predicted property values
match the eight that were requested?}
\label{fig:r2two}
\end{figure}

\subsection{The paper's definition, precisely}

The released notebook makes it unambiguous. Inside the design loop:

\begin{lstlisting}[style=py]
r2 = r2_score(req, prop)   # req: the 8 requested values
                           # prop: the 8 predicted values
\end{lstlisting}

Both arguments have length 8. So $n = 8$, and the eight ``observations'' are the
eight \emph{properties} of a single sequence.

\begin{keybox}{What the paper's \Rsq{} actually measures}
It is a \term{self-consistency} check on a round trip:

\begin{center}
\begin{tikzpicture}[font=\small,
  n/.style={draw=greyc!70,fill=boxbg,rounded corners=2pt,inner sep=5pt,
            text width=2.6cm,align=center,minimum height=0.95cm},
  ar/.style={-{Stealth[length=2.4mm]},thick,greyc!80},
]
\node[n,fill=papercol!10,draw=papercol!60] (t) {target\\ 8 numbers};
\node[n,right=1.5cm of t] (s) {designed\\ sequence};
\node[n,right=1.5cm of s] (p) {predicted\\ 8 numbers};
\draw[ar] (t) -- node[above,font=\scriptsize]{inverse} (s);
\draw[ar] (s) -- node[above,font=\scriptsize]{forward} (p);
\draw[ar,ourscol,dashed] (p.south) to[bend left=32]
  node[below,font=\scriptsize,text=ourscol]{compare: this is the $R^2$} (t.south);
\end{tikzpicture}
\end{center}

It asks: when the model designs a sequence for a target, and is then asked what
that sequence's properties are, does it agree with itself? That is a
meaningful thing to measure, and the paper describes it as such. It is
\emph{not} a measure of whether the sequence really has those properties, and
it is not a measure of generalisation to unseen sequences.
\end{keybox}

\subsection{Two consequences that bite}

\paragraph{The denominator is tiny, and varies between targets.}
$\mathrm{SS}_{\text{tot}}$ is the spread of the \emph{eight target values}. If a
target is $[0.25, 0.252, 0.20, 0.151, 0.31, 0.203, 0.293, 0.265]$ --- all
clustered around $0.25$ --- then the denominator is small, and a modest absolute
error produces a dramatic \Rsq{}. Two targets with the same absolute accuracy
can score very differently purely because one is more spread out. I report this
spread alongside every within-vector \Rsq{} in the project, under the name
\code{target\_spread}.

\paragraph{Half the entries are standard deviations.}
Four of the eight numbers are property values; four are measurement standard
deviations. The SDs vary less across the dataset, so they are easier to predict.
A model that nails the four SDs and misses the four properties still scores
respectably. Throughout the project I therefore report the four mechanical
slots separately, via an index list \code{MECHANICAL\_IDX = [0, 2, 4, 6]}.

\begin{warnbox}{Argument order is not symmetric}
$R^2(a, b) \neq R^2(b, a)$, because the denominator uses the variance of the
\emph{first} argument only. The released notebook uses \code{r2\_score(req,
prop)} --- target first --- in the loop that produces the published numbers, and
\code{r2\_score(c\_res, c\_GT)} --- predictions first --- in a separate helper
that does not appear to feed any published figure. I mention it only because
somebody reusing that helper would be computing a different quantity. This
project always uses (target, prediction), and has a test asserting the two
differ so the convention cannot drift.
\end{warnbox}

\section{Best-of-\texorpdfstring{$N$}{N}: why a maximum is not a measurement}
\label{sec:bestofn-concept}

This is the single most important statistical idea in the project.

Suppose a procedure produces candidates with some spread of quality. Generate
$N$ of them and keep the best. The expected quality of what you keep
\emph{increases with $N$}, without anything about the procedure improving.

\begin{figure}[htbp]
\centering
\begin{tikzpicture}[font=\small]
\begin{axis}[
  width=12.5cm,height=5.6cm,
  xlabel={quality of a single candidate (e.g.\ $R^2$)},
  ylabel={how often},
  axis lines=left, ytick=\empty,
  xmin=-3.2,xmax=1.15, ymin=0, ymax=1.15,
  xlabel style={font=\scriptsize}, ylabel style={font=\scriptsize},
  xtick={-3,-2,-1,0,1}, xticklabel style={font=\scriptsize},
  clip=false,
]
\addplot[draw=papercol,fill=papercol!18,domain=-3.2:1.0,samples=200]
  {exp(-((x+0.9)^2)/(2*0.55^2))} \closedcycle;

\draw[ourscol,line width=1.1pt,dashed] (axis cs:0.72,0) -- (axis cs:0.72,1.02);
\node[anchor=south,font=\scriptsize,text=ourscol,align=center]
  at (axis cs:0.72,1.02) {the one you keep\\ after $N$ draws};

\draw[greyc,line width=0.9pt,dotted] (axis cs:-0.9,0) -- (axis cs:-0.9,0.92);
\node[anchor=south,font=\scriptsize,text=greyc] at (axis cs:-1.5,0.92)
  {typical draw};

\draw[decorate,decoration={brace,amplitude=4pt},greyc!80]
  (axis cs:-0.85,-0.16) -- (axis cs:0.70,-0.16)
  node[midway,below=4pt,font=\scriptsize,text=greyc]
  {this gap is produced by \emph{searching}, not by the model};
\end{axis}
\end{tikzpicture}
\caption[Best-of-N]{The best-of-$N$ effect. The shaded curve is the quality of
a single candidate --- mostly poor. Draw many and keep the maximum, and you land
far out in the right tail. Report only the maximum, without $N$, and a reader
cannot tell how much of the number came from the model and how much from the
search.}
\label{fig:bestofn-concept}
\end{figure}

\begin{keybox}{The order statistic}
If a single candidate's quality has cumulative distribution $F$, then the
maximum of $N$ independent draws has distribution
\[
  P(\max \le x) = F(x)^N ,
\]
which pushes steadily rightward as $N$ grows. The expected maximum increases
without bound in $N$ (until it saturates at the distribution's upper limit).
This is a fact about sampling, not about the sampler's competence.
\end{keybox}

\begin{warnbox}{Why this matters here}
The paper reports its \Rsq{} values as ``the results with the highest \Rsq{}
values are selected from a set of sampling attempts''. The size of that set
appears nowhere in the paper. The released notebook shows it: 32 samples per
step $\times$ 64 repeats $=$ up to \textbf{2{,}048 candidates}, then
\code{np.argmax}.

There is nothing wrong with the procedure. If generation is cheap and screening
is cheaper, generating thousands of candidates and keeping the best is exactly
what a practitioner should do. The issue is purely one of reporting: without
$N$, the number cannot be compared against another method's single-shot result,
and a reader cannot estimate the compute needed to reach it.

Quantifying this became Extension~1 (Section~\ref{sec:bestofn}).
\end{warnbox}

\section{In-sample versus held-out, and why it matters here}

A last piece of vocabulary, needed for Chapter~\ref{ch:extensions}.

\begin{description}[leftmargin=0cm,style=nextline,itemsep=4pt]
\item[In-sample.] Performance measured on the data the model was trained on. It
  is an upper bound. A model that simply memorised a lookup table scores
  perfectly in-sample and has learned nothing.

\item[Held-out (out-of-sample).] Performance on data withheld from training.
  This is what tells you whether anything general was learned.
\end{description}

Since SilkomeGPT was fine-tuned on all 1{,}033 pairs, every property-prediction
number measurable on those pairs --- the paper's and mine --- is in-sample. This
does not contradict anything the paper claims, because it does not claim
held-out accuracy. But it means the two possibilities, \emph{learning} and
\emph{memorising}, cannot be told apart from the published figures alone.
Distinguishing them is Extension~4 (Section~\ref{sec:memorisation}), and the
answer is not the comfortable one.
